\section{Introduction}
\label{sec:introduction}

Conformance checking is useful because a process model and an event log describe different things. The model states which executions are allowed; the log records what was observed. A difference between them may reveal an exception, an incomplete recording, or a model that needs revision. An alignment makes this difference concrete by relating each recorded event to an execution of the model and identifying the extra or missing behavior needed to reconcile the two~\cite{Aalst2016,Adriansyah2011,Carmona2018}. Its value lies in both a deviation cost and an explanation that can be replayed on the model. Figure~\ref{fig:outline} outlines how the paper develops and evaluates this approach.

\begin{figure*}[!t]
\centering
\begin{tikzpicture}[x=1cm,y=1cm,
  stage/.style={text width=26mm,minimum height=19mm,inner sep=5pt},
  result/.style={text width=53mm,minimum height=20mm,inner sep=5pt}]

\node[paneltitle] at (0,3.6) {TRAIN ONCE};
\node[box,stage] (data) at (1.5,2.35)
  {\textbf{Synthetic data}\\Petri nets + traces\\exact alignments};
\node[learnbox,stage] (training) at (5,2.35)
  {\textbf{Train scorer}\\learn move\\preferences};
\draw[flow] (data) -- (training);
\node[annotation,anchor=west,text width=58mm] at (7.15,2.35)
  {One checkpoint reused\\across process models\\and activity vocabularies.};

\node[paneltitle] at (0,.85) {ALIGN EACH NEW INPUT};
\node[box,stage] (input) at (1.5,-.3)
  {\textbf{New input}\\Petri net + trace\\initial/final\\markings};
\node[learnbox,stage] (scorer) at (5,-.3)
  {\textbf{Score moves}\\encode graph\\and trace\\rank alternatives};
\node[learnbox,stage] (decoder) at (8.5,-.3)
  {\textbf{Build candidate}\\track tokens\\follow firing rules};
\node[verifybox,stage] (replay) at (12,-.3)
  {\textbf{Replay check}\\validate execution\\recompute cost};
\draw[flow,dashed,draw=learnblue] (training) --
  node[right,font=\small,text=learnblue] {fixed parameters} (scorer);
\draw[flow] (input) -- (scorer);
\draw[flow] (scorer) -- (decoder);
\draw[flow] (decoder) -- (replay);

\node[verifybox,result] (verified) at (10.5,-3.45)
  {\textbf{Verified alignment}\\executable explanation\\cost upper bound};
\draw[flow] (replay.south) -- (12,-2.25) -|
  node[pos=.25,above,font=\small] {replay passes} (verified.north);

\node[verifybox,result] (exact) at (3.25,-3.45)
  {\textbf{Optional exact certification}\\solve the original input\\retain or replace the candidate};
\draw[flow,dashed] (input.south) -- (1.5,-2.25) -| (exact.north);
\draw[flow,dashed] (replay.south) -- (12,-1.75) -- (5.6,-1.75) --
  node[right,font=\small,text=muted] {candidate} ([xshift=23.5mm]exact.north);
\node[annotation,anchor=north,text width=59mm,align=center] at (3.25,-4.65)
  {Optimality requires\\successful exact search.};
\node[annotation,anchor=north,text width=57mm,align=center] at (10.5,-4.65)
  {Replay failure: no verified\\candidate or cost bound.};
\end{tikzpicture}

\caption{Outline of the paper. The argument moves from concepts and related work to the proposed method, its training data, and its implementation, then evaluates the results and their limitations. Arrows indicate the reading order.}
\label{fig:outline}
\end{figure*}

For example, suppose that two transitions have the same activity label, but lead to different branches. The next event alone cannot identify the appropriate transition. Later events may reveal which branch gives the best explanation. An exact aligner explores alternatives and returns a minimum cost explanation. This is an attractive contract, but exploration becomes expensive when many choices interact through concurrency, repetition, or invisible routing. Even when most traces are inexpensive, a few difficult variants can delay a complete analysis. Faster first results could make conformance checking more responsive, provided their meaning is clear.

This paper investigates learned guidance for that first result. A neural network can learn recurring relationships between process structure, observed events, and good alignment choices. Once trained, it can assign scores to moves on a new net--trace pair. The difficulty is that a high score does not prove that a move is enabled, that a complete execution will reach the final marking, or that the resulting cost is minimal. Consequently, replacing formal search with a predictor would change the meaning of the answer unless those obligations were handled explicitly.

LARA, short for \emph{Learned Alignment with Replay Assurance}, assigns these obligations to different components. The neural network proposes preferences. A decoder follows Petri net firing rules while constructing a candidate. An independent replay routine checks the complete candidate and computes its cost. An optional exact aligner then determines whether that candidate is optimal and replaces it if necessary. The basic principle is simple: a prediction may influence which explanation is tried first, while process semantics determines which claims can accompany it.

The reuse question is central. Training a separate neural network for every process would create a substantial preparation burden and would weaken the motivation for machine learning in the first place. We therefore train one scorer on abstract synthetic process structures and reuse its parameters on other inputs. Activity names identify equal labels; the neural network does not interpret their business meaning. The study follows the idea of pretraining a reusable model, but its scope is narrow: control flow alignment under fixed unit deviation costs. We use the term \emph{foundation model approach} in this limited sense and distinguish it from the broad task coverage often associated with foundation models~\cite{Bommasani2021,Rizk2022}.

The empirical conclusion is similarly specific. Learned scores improve the same constrained decoder from 86.9\% to 91.0\% optimal candidates on the synthetic test split, with the strongest gain for duplicate label choices. All evaluated candidates pass replay. However, exact search is faster on small inputs, and several established approximation methods produce better candidates overall. The present certified mode also runs exact search on every trace; the proportion of candidates that are retained is therefore a quality measure, not a reduction in exact computation. These findings motivate a careful account of where machine learning contributes and where the decoder or exact backend does most of the work.

The paper follows the outline in Figure~\ref{fig:outline}. After introducing the concepts and related work, it explains the method and its synthetic training procedure, describes the implementation, and evaluates candidate quality, the contribution of machine learning, and transfer. The discussion connects these findings to their practical limits.
